\section{Quantitative Approximation by the Mean-Field Limit for Low-Rank Networks}
\label{sec:quantitative}

This section provides a rigorous quantitative bound on the approximation error between finite-width low-rank networks and their mean-field limit. The key result shows that the approximation error scales as $O(1/\sqrt{n_{\min}} + \sqrt{\epsilon})$ with explicit constants, where $n_{\min}$ is the minimum width across layers and $\epsilon$ is the learning rate step size. This bound holds for any $n_1$ and $n_2$, independent of the data dimension $d$, similar to the full-width case \cite{nguyen2020global,araujo2019mean}. The bound suggests that widths $n_1, n_2 \approx 1000$ are typically sufficient to observe mean-field behaviors, as empirically validated in \cite{nguyen2019meanfield} for high-dimensional real-world data.

\subsection{Main Result: Approximation by the MF Limit}

\begin{assumption}[Initialization for Low-Rank Networks]
\label{assump:quant-init}
We assume that ${\rm ess\text{-}sup}\max_{1\le k\le r}|w_1^0(C_1,k)| \le K$ and ${\rm ess\text{-}sup}|w_2^0(C_2)| \le K$, where $w_1^0$ and $w_2^0$ are the initial weights as described in the low-rank architecture setup.
\end{assumption}

\begin{theorem}[Finite-width approximation error bound]
\label{thm:quantitative-lowrank}
Given a family $\mathsf{Init}$ of initialization laws and a tuple $\{n_1, n_2\}$ that is in the index set of $\mathsf{Init}$, perform the coupling procedure for the low-rank architecture as described in Section~\ref{sec:setup}. Fix a terminal time $T \in \epsilon\mathbb{N}_{\ge 0}$. Under Assumptions~\ref{assump:regularity}, \ref{assump:quant-init}, and the low-rank structure with mixing matrix $L$ satisfying $\|L\|_{\infty,1} \le \textcolor{red}{rK}$, for $\epsilon \le 1$, we have with probability at least $1-2\delta$,
\begin{align*}
\mathscr{D}_T(W, \mathbf{W}) &\le \mathbf{e^{K_T\textcolor{red}{(1+rK)}}}\left(\frac{1}{\sqrt{n_{\min}}} + \sqrt{\epsilon}\right)\log^{1/2}\left(\frac{3(T+1)n_{\max}^2}{\delta} + e\right) \\
&\equiv \mathsf{err}_{\delta,T}(\epsilon, n_1, n_2, r),
\end{align*}
in which $n_{\min} = \min\{n_1, n_2\}$, $n_{\max} = \max\{n_1, n_2\}$, $K_T = K(1+T^K)$, and the factor $\textcolor{red}{(1+rK)}$ accounts for the low-rank structure through $\|L\|_{\infty,1} \le \textcolor{red}{rK}$.
\end{theorem}

The theorem gives a connection between $\mathbf{W}(\lfloor t/\epsilon\rfloor)$ (the discrete-time finite-width network) and $W(t)$ (the continuous-time mean-field limit). This result extends the quantitative approximation theorem from the full-width case \cite{nguyen2020global} to the low-rank setting, with \textbf{the key difference being the multiplicative factor $\textcolor{red}{(1+rK)}$ in the exponential constant}, which accounts for the low-rank mixing matrix structure.

\begin{corollary}[Test function approximation quality]
\label{cor:gradient-descent-quality-lowrank}
Under the same setting as Theorem~\ref{thm:quantitative-lowrank}, consider any test function $\psi:\mathbb{R}\times\mathbb{R}\to\mathbb{R}$ which is $K$-Lipschitz in the second variable uniformly in the first variable (an example of $\psi$ is the loss $\mathscr{L}$). For any $\delta>0$, with probability at least $1-3\delta$,
\[
\sup_{t\leq T}\left|\mathbb{E}_{Z}\left[\psi\left(Y,\hat{\mathbf{y}}\left(X;\mathbf{W}\left(\left\lfloor t/\epsilon\right\rfloor \right)\right)\right)\right]-\mathbb{E}_{Z}\left[\psi\left(Y,\hat{y}\left(X;W\left(t\right)\right)\right)\right]\right|\leq e^{K_T\textcolor{red}{(1+rK)}}\mathsf{err}_{\delta,T}\left(\epsilon,n_{1},n_{2},r\right),
\]
where $\hat{\mathbf{y}}(X; \mathbf{W})$ and $\hat{y}(X; W)$ denote the outputs of the finite-width and mean-field networks respectively.
\end{corollary}

\begin{proof}[Proof sketch]
The proof follows the same structure as in the full-width case \cite{nguyen2020global}. Since $\psi$ is $K$-Lipschitz in the second variable, we have:
\[
|\psi(Y, \hat{\mathbf{y}}(X; \mathbf{W})) - \psi(Y, \hat{y}(X; W))| \le K|\hat{\mathbf{y}}(X; \mathbf{W}) - \hat{y}(X; W)|.
\]
The difference $|\hat{\mathbf{y}}(X; \mathbf{W}) - \hat{y}(X; W)|$ can be bounded by the distance $\mathscr{D}_T(W, \mathbf{W})$ using the low-rank structure of the network. Specifically, for the low-rank architecture, the output difference involves:
\begin{align*}
|\hat{\mathbf{y}}(X; \mathbf{W}) - \hat{y}(X; W)| &\le \sum_{k=1}^r |L_{C_2,k}|\left|\frac{1}{n_1}\sum_{j_1=1}^{n_1}\mathbf{w}_1(\lfloor t/\epsilon\rfloor, j_1, k)\varphi_1(\langle f_1(C_1(j_1)), X\rangle)\right.\\
&\quad\left. - \mathbb{E}_{C_1}[w_1(t, C_1, k)\varphi_1(\langle f_1(C_1), X\rangle)]\right| + \text{(similar terms for $w_2$)}\\
&\le \textcolor{red}{\mathbf{rK}}\mathscr{D}_T(W, \mathbf{W}) + K\mathscr{D}_T(W, \mathbf{W}) = \textcolor{red}{\mathbf{K(1+rK)}}\mathscr{D}_T(W, \mathbf{W}),
\end{align*}
where we used $\|L\|_{\infty,1} \le \textcolor{red}{rK}$. Taking expectation over $Z$ and applying Theorem~\ref{thm:quantitative-lowrank} completes the proof.
\end{proof}

\paragraph*{Comparison with the Full-Width Case.}
\textbf{The main difference between our low-rank result and the full-width case \cite{nguyen2020global} is the appearance of the factor $\textcolor{red}{(1+rK)}$ in the exponential constant.} In the full-width case, the bound is:
\[
\mathscr{D}_T(W, \mathbf{W}) \le e^{K_T}\left(\frac{1}{\sqrt{n_{\min}}} + \sqrt{\epsilon}\right)\log^{1/2}\left(\frac{3(T+1)n_{\max}^2}{\delta} + e\right),
\]
whereas in our low-rank case, we have:
\[
\mathscr{D}_T(W, \mathbf{W}) \le \mathbf{e^{K_T\textcolor{red}{(1+rK)}}}\left(\frac{1}{\sqrt{n_{\min}}} + \sqrt{\epsilon}\right)\log^{1/2}\left(\frac{3(T+1)n_{\max}^2}{\delta} + e\right).
\]

The factor $\textcolor{red}{(1+rK)}$ arises from:
\begin{itemize}
\item The low-rank structure of $H_2$, which involves a sum over $r$ channels: $H_2(t, c_2; X, W) = \sum_{k=1}^r L_{c_2,k}\,m_k(t; X, W)$.
\item The mixing matrix bounds: $\|L\|_{\infty,1} = \sup_{c_2}\sum_{k=1}^r |L_{c_2,k}| \le \textcolor{red}{rK}$.
\item The backpropagated signal $B_k$ which aggregates over $n_2$ neurons with mixing coefficients $L_{C_2,k}$.
\end{itemize}

However, the fundamental structure of the proof remains the same: we still decompose the error into particle coupling and gradient descent discretization, and the scaling with $n_{\min}$ and $\epsilon$ is identical. The low-rank structure only affects the exponential constant, not the polynomial scaling. This suggests that the mean-field approximation quality is preserved under low-rank constraints, with the trade-off being a potentially larger (but still finite) exponential constant.

\begin{remark}[Deterministic mixing matrix and hierarchical learning]
The mixing matrix $L$ can be chosen deterministically (e.g., on a grid) rather than randomly, since the proof only requires the boundedness condition $\|L\|_{\infty,1} \le \textcolor{red}{rK}$. There is no advantage to maximizing the entries of $L$; fixing $L$ deterministically with appropriate structure can enable hierarchical learning, where different channels $k$ specialize to different frequency components or scales through the backpropagated signal $B_k$. This design choice allows for structured learning dynamics while maintaining the same theoretical guarantees.
\end{remark}

\begin{remark}[Low-rank framework enables i.i.d. initialization for arbitrary depth]
\label{rem:iid-init-lowrank}
A key theoretical advantage of our low-rank framework is that it enables standard i.i.d. initialization for networks of arbitrary depth $L$, whereas full-width networks require correlated initialization for $L \geq 3$.

In the full-width case, by Corollary 29 of \cite{nguyen2023rigorousframeworkmeanfield}, neurons at intermediate layers collapse to the same function of the input under i.i.d. initialization for $L \geq 3$, and therefore are not expected to span the space of functions of the input $L^2(\mathcal{P}_X)$. In other words, these intermediate layers become a bottleneck that hinders universal approximation in the context of more than three layers and general data distributions. To attain meaningful training, this suggests a departure from i.i.d. initializations, requiring a correlated initialization scheme that introduces correlation among weights across layers.

\emph{Our low-rank framework with frozen random features eliminates this requirement entirely.} The frozen first-layer features $f_1(C_1)$ automatically maintain full support $\text{supp}(f_1(C_1)) = \mathbb{R}^d$ throughout training, independent of the target function. This means:
\begin{itemize}
\item \textbf{Standard i.i.d. initialization suffices}: We can initialize $w_1^0(C_1, k)$ and $w_2^0(C_2)$ independently from standard distributions (e.g., Gaussian), without requiring any correlation structure.
\item \textbf{Function-independent}: The initialization does not depend on the target function to approximate, unlike the correlated initialization scheme needed for full-width networks.
\item \textbf{Arbitrary depth}: Global convergence holds for any depth $L$, not just $L=3$, with the same simple initialization scheme.
\end{itemize}
This represents a significant theoretical improvement: while full-width networks require carefully designed, function-dependent correlated initialization for $L \geq 3$, our low-rank framework achieves global convergence for arbitrary depth with standard i.i.d. initialization. The frozen random features act as a stable basis that prevents the neuron collapse phenomenon observed in full-width networks.
\end{remark}

\begin{remark}[Generalization to $L$ layers: recursive pattern]
\label{rem:generalization-L}
The proof strategy generalizes naturally to arbitrary depth $L$ because \textbf{the same three-step pattern repeats at each intermediate layer}:

\begin{enumerate}
\item \textbf{Forward decomposition}: For each $H_i$ with $i = 2, \ldots, L$, we decompose the difference $q_i = \tilde{H}_i - H_i$ using the same structure:
\[
|q_i(t, x, j_i, c_i)| \le \sum_{k=1}^r |L^{(i)}_{C_i(j_i),k}|(Q_{i,1,k} + Q_{i,2,k}) \le rK\,\max_{1 \le k \le r}(Q_{i,1,k} + Q_{i,2,k}),
\]
where $Q_{i,1,k}$ bounds weight differences and $Q_{i,2,k}$ bounds concentration errors. This pattern is identical for all layers.

\item \textbf{Backward decomposition}: For each backpropagated signal $B_k^{(i)}$ with $i = 2, \ldots, L-1$, we decompose:
\[
|q_{B,k}^{(i)}(t, x)| \le Q_{B,1,k}^{(i)}(x) + Q_{B,2,k}^{(i)}(x),
\]
where $Q_{B,1,k}^{(i)}$ bounds activation/weight differences and $Q_{B,2,k}^{(i)}$ bounds concentration errors. This pattern is identical for all intermediate layers.

\item \textbf{Concentration and Gronwall}: The concentration bounds and Gronwall argument proceed identically at each layer, with union bounds over $r$ channels and factors of $rK$ from $\|L^{(i)}\|_{\infty,1} \le rK$.
\end{enumerate}

\textbf{Key insight:} With bounded mixing matrices $\|L^{(i)}\|_{\infty,1} \le rK$, bounded activations (e.g., ReLU), and bounded data $|X| \le K$, the exponential constant accumulates multiplicatively as $(1+rK)^{L-2}$ rather than exponentially in $L$. The frozen first layer eliminates recursion at the base case, ensuring the pattern starts cleanly. Each layer applies the same bounded operations, so depth $L$ only enters polynomially in the exponential constant, not exponentially. This is why the proof works for arbitrary depth: \textbf{the low-rank structure with bounded mixing matrices ensures that each layer contributes at most a factor of $(1+rK)$ to the exponential constant, and these factors multiply rather than compound exponentially.}

\textbf{Crucial point:} The frozen first layer is not just a simplification—it's the key mechanism that prevents divergence. Since $H_1(t, c_1; X, W) = \langle f_1(c_1), X\rangle$ is independent of $W(t)$, it cannot diverge and maintains full support $\text{supp}(f_1(C_1)) = \mathbb{R}^d$ trivially. This stable base allows us to apply the same bounded procedure recursively at each subsequent layer, with all errors bounded in terms of the global distance $\mathscr{D}_t(W, \tilde{W})$ rather than layer-wise errors that would compound exponentially. Without the frozen first layer, we would need correlated initialization (as in the full-width case) and would be limited to $L=3$ layers.
\end{remark}

\subsection{Particle ODEs for Low-Rank Networks}

We construct auxiliary trajectories, which we call the \textit{particle ODEs} for the low-rank case. These are continuous-time trajectories of finitely many neurons, averaged over the data distribution, adapted to the low-rank structure:
\begin{align*}
\frac{\partial}{\partial t}\tilde{w}_2(t, j_2) &= -\xi_2(t)\mathbb{E}_Z\left[d_L(Z; \tilde{W}(t))\,\varphi_2(\tilde{H}_2(t, j_2; X, \tilde{W}(t)))\right],\\
\frac{\partial}{\partial t}\tilde{w}_1(t, j_1, k) &= -\xi_1(t)\mathbb{E}_Z\left[d_L(Z; \tilde{W}(t))\,\varphi_1(\langle f_1(C_1(j_1)), X\rangle)\,\tilde{B}_k(t; X, \tilde{W}(t))\right],
\end{align*}
where $j_1 = 1, \ldots, n_1$, $j_2 = 1, \ldots, n_2$, $k = 1, \ldots, r$, $\tilde{W}(t) = (\tilde{w}_1(t, \cdot, \cdot), \tilde{w}_2(t, \cdot))$, and $t \in \mathbb{R}_{\ge 0}$. 

The second-layer output for the particle ODEs is:
\[
\tilde{H}_2(t, j_2; X, \tilde{W}) = \sum_{k=1}^r L_{C_2(j_2),k}\,\frac{1}{n_1}\sum_{j_1=1}^{n_1} \tilde{w}_1(t, j_1, k)\,\varphi_1(\langle f_1(C_1(j_1)), X\rangle),
\]
and the backpropagated signal is:
\[
\tilde{B}_k(t; X, \tilde{W}) = \frac{1}{n_2}\sum_{j_2=1}^{n_2} L_{C_2(j_2),k}\,\varphi_2'(\tilde{H}_2(t, j_2; X, \tilde{W}))\,\tilde{w}_2(t, j_2).
\]

We specify the initialization $\tilde{W}(0)$: $\tilde{w}_1(0, j_1, k) = w_1^0(C_1(j_1), k)$ and $\tilde{w}_2(0, j_2) = w_2^0(C_2(j_2))$. That is, it shares the same initialization with the neural network $\mathbf{W}(0)$, and hence is coupled with the neural network and the MF ODEs.

The existence and uniqueness of the solution to the particle ODEs follows from the same proof as in Theorem~\ref{thm:wellposed}, adapted to account for the low-rank structure. We equip $\tilde{W}(t)$ with the norm:
\[
\interleave\tilde{W}\interleave_T = \max\left\{\max_{j_1 \le n_1, 1 \le k \le r}\sup_{t \le T}|\tilde{w}_1(t, j_1, k)|, \max_{j_2 \le n_2}\sup_{t \le T}|\tilde{w}_2(t, j_2)|\right\}.
\]

We define the distance measures:
\begin{align*}
\mathscr{D}_T(W, \tilde{W}) &= \sup\Big\{|w_1(t, C_1(j_1), k) - \tilde{w}_1(t, C_1(j_1), k)|, |w_2(t, C_2(j_2)) - \tilde{w}_2(t, C_2(j_2))|:\\
&\qquad t \le T, j_1 \le n_1, j_2 \le n_2, 1 \le k \le r\Big\},\\
\mathscr{D}_T(\tilde{W}, \mathbf{W}) &= \sup\Big\{|\mathbf{w}_1(\lfloor t/\epsilon\rfloor, j_1, k) - \tilde{w}_1(t, C_1(j_1), k)|,\\
&\qquad |\mathbf{w}_2(\lfloor t/\epsilon\rfloor, j_2) - \tilde{w}_2(t, C_2(j_2))|: t \le T, j_1 \le n_1, j_2 \le n_2, 1 \le k \le r\Big\}.
\end{align*}

\subsection{Key Lemmas}

\begin{lemma}[Particle coupling bound]
\label{lem:particle-lowrank}
Under the same setting as Theorem~\ref{thm:quantitative-lowrank}, for any $\delta > 0$, with probability at least $1-\delta$,
\[
\mathscr{D}_T(W, \tilde{W}) \le \frac{1}{\sqrt{n_{\min}}}\log^{1/2}\left(\frac{3(T+1)n_{\max}^2}{\delta} + e\right)\mathbf{e^{K_T\textcolor{red}{(1+rK)}}},
\]
in which $n_{\min} = \min\{n_1, n_2\}$, $n_{\max} = \max\{n_1, n_2\}$, and $K_T = K(1+T^K)$.
\end{lemma}

\begin{lemma}[Gradient descent discretization error]
\label{lem:gradient-lowrank}
Under the same setting as Theorem~\ref{thm:quantitative-lowrank}, for any $\delta > 0$ and $\epsilon \le 1$, with probability at least $1-\delta$,
\[
\mathscr{D}_T(\tilde{W}, \mathbf{W}) \le \sqrt{\epsilon\log\left(\frac{\mathbf{2n_1n_2r}}{\delta} + e\right)}\mathbf{e^{K_T\textcolor{red}{(1+rK)}}},
\]
in which $K_T = K(1+T^K)$.
\end{lemma}

\begin{proof}[Proof of Theorem~\ref{thm:quantitative-lowrank}]
Using the triangle inequality:
\[
\mathscr{D}_T(W, \mathbf{W}) \le \mathscr{D}_T(W, \tilde{W}) + \mathscr{D}_T(\tilde{W}, \mathbf{W}),
\]
the result follows immediately from Lemmas~\ref{lem:particle-lowrank} and \ref{lem:gradient-lowrank}, noting that the $\log(2n_1n_2r/\delta)$ term can be absorbed into the $\log(3(T+1)n_{\max}^2/\delta + e)$ term up to constants.
\end{proof}

\subsection{Proof of Lemma~\ref{lem:particle-lowrank}}

\begin{proof}
In the following, let $K_t$ denote a generic positive constant that may change from line to line and takes the form $K_t = K(1+t^K)$, such that $K_t \ge 1$ and $K_t \le K_T$ for all $t \le T$. We first note that at initialization, $\mathscr{D}_0(W, \tilde{W}) = 0$. Since $\interleave W\interleave_0 \le K$, we have $\interleave W\interleave_T \le K_T$ by the a priori bounds. Furthermore, it is easy to see that $\interleave\tilde{W}\interleave_0 \le \interleave W\interleave_0 \le K$ almost surely. By the same argument, $\interleave\tilde{W}\interleave_T \le K_T$ almost surely.

We decompose the proof into several steps.

\paragraph*{Step 1: Main bound with low-rank structure.}

Let us define, for brevity, the differences specific to the low-rank architecture:
\begin{align*}
q_2(t, x, j_2, c_2) &= \tilde{H}_2(t, j_2; x, \tilde{W}(t)) - H_2(t, c_2; x, W(t)),\\
q_{B,k}(t, x) &= \tilde{B}_k(t; x, \tilde{W}(t)) - B_k(t; x, W(t)),
\end{align*}
where $H_2(t, c_2; x, W) = \sum_{k=1}^r L_{c_2,k}\,\mathbb{E}_{C_1}[w_1(t, C_1, k)\,\varphi_1(\langle f_1(C_1), x\rangle)]$ and $B_k(t; x, W) = \mathbb{E}_{C_2}[L_{C_2,k}\,\varphi_2'(H_2(t, C_2; x, W))\,w_2(t, C_2)]$.

Consider $t \ge 0$. We first bound the difference in the updates between $W$ and $\tilde{W}$.

\subparagraph*{Bound for $w_2$ and $\tilde{w}_2$:}
By Assumption~\ref{assump:regularity} and the definition of the ODEs:
\begin{align*}
&\left|\frac{\partial}{\partial t}\tilde{w}_2(t, j_2) - \frac{\partial}{\partial t}w_2(t, C_2(j_2))\right|\\
&\quad = \left|\xi_2(t)\mathbb{E}_Z\left[d_L(Z; \tilde{W}(t))\,\varphi_2(\tilde{H}_2(t, j_2; X, \tilde{W}(t))) - d_L(Z; W(t))\,\varphi_2(H_2(t, C_2(j_2); X, W(t)))\right]\right|\\
&\quad \le K\mathbb{E}_Z\left[|d_L(Z; \tilde{W}(t)) - d_L(Z; W(t))|\,|\varphi_2(\tilde{H}_2(t, j_2; X, \tilde{W}(t)))|\right]\\
&\quad\quad + K\mathbb{E}_Z\left[|d_L(Z; W(t))|\,|\varphi_2(\tilde{H}_2(t, j_2; X, \tilde{W}(t))) - \varphi_2(H_2(t, C_2(j_2); X, W(t)))|\right]\\
&\quad \le K_t\mathbb{E}_Z\left[|q_2(t, X, j_2, C_2(j_2))|\right] + K_t\mathscr{D}_t(W, \tilde{W}),
\end{align*}
where we used that $d_L$ is Lipschitz in $W$, $\varphi_2$ is Lipschitz, and $H_2$ differences are controlled.

\subparagraph*{Bound for $w_1$ and $\tilde{w}_1$:}
For the low-rank case, we have $k = 1, \ldots, r$ channels. By the definition of the ODEs:
\begin{align*}
&\left|\frac{\partial}{\partial t}\tilde{w}_1(t, j_1, k) - \frac{\partial}{\partial t}w_1(t, C_1(j_1), k)\right|\\
&\quad = \left|\xi_1(t)\mathbb{E}_Z\left[d_L(Z; \tilde{W}(t))\,\varphi_1(\langle f_1(C_1(j_1)), X\rangle)\,\tilde{B}_k(t; X, \tilde{W}(t))\right.\right.\\
&\quad\quad\left.\left. - d_L(Z; W(t))\,\varphi_1(\langle f_1(C_1(j_1)), X\rangle)\,B_k(t; X, W(t))\right]\right|\\
&\quad \le K\mathbb{E}_Z\left[|d_L(Z; \tilde{W}(t)) - d_L(Z; W(t))|\,|\tilde{B}_k(t; X, \tilde{W}(t))|\right]\\
&\quad\quad + K\mathbb{E}_Z\left[|d_L(Z; W(t))|\,|q_{B,k}(t, X)|\right]\\
&\quad \le K_t\mathbb{E}_Z\left[|q_{B,k}(t, X)|\right] + K_t\mathscr{D}_t(W, \tilde{W}),
\end{align*}
where the expectation over $j_2$ in $\tilde{B}_k$ will be handled via concentration.

\paragraph*{Step 2: Decomposition of $q_2$ with low-rank structure.}

\textbf{The key difference from the full-width case is that $H_2$ involves a sum over $r$ channels.} We decompose:
\begin{align*}
|q_2(t, x, j_2, c_2)| &= \left|\textcolor{red}{\sum_{k=1}^r L_{C_2(j_2),k}}\left(\frac{1}{n_1}\sum_{j_1=1}^{n_1}\tilde{w}_1(t, j_1, k)\,\varphi_1(\langle f_1(C_1(j_1)), x\rangle)\right.\right.\\
&\quad\left.\left. - \mathbb{E}_{C_1}[w_1(t, C_1, k)\,\varphi_1(\langle f_1(C_1), x\rangle)]\right)\right|\\
&\le \textcolor{red}{\sum_{k=1}^r |L_{C_2(j_2),k}|}\left|\frac{1}{n_1}\sum_{j_1=1}^{n_1}\tilde{w}_1(t, j_1, k)\,\varphi_1(\langle f_1(C_1(j_1)), x\rangle)\right.\\
&\quad\left. - \mathbb{E}_{C_1}[w_1(t, C_1, k)\,\varphi_1(\langle f_1(C_1), x\rangle)]\right|.
\end{align*}

We further decompose each term in the sum:
\begin{align*}
&\left|\frac{1}{n_1}\sum_{j_1=1}^{n_1}\tilde{w}_1(t, j_1, k)\,\varphi_1(\langle f_1(C_1(j_1)), x\rangle) - \mathbb{E}_{C_1}[w_1(t, C_1, k)\,\varphi_1(\langle f_1(C_1), x\rangle)]\right|\\
&\quad \le \max_{j_1 \le n_1}\left|\tilde{w}_1(t, j_1, k)\,\varphi_1(\langle f_1(C_1(j_1)), x\rangle) - w_1(t, C_1(j_1), k)\,\varphi_1(\langle f_1(C_1(j_1)), x\rangle)\right|\\
&\quad\quad + \left|\frac{1}{n_1}\sum_{j_1=1}^{n_1}w_1(t, C_1(j_1), k)\,\varphi_1(\langle f_1(C_1(j_1)), x\rangle) - \mathbb{E}_{C_1}[w_1(t, C_1, k)\,\varphi_1(\langle f_1(C_1), x\rangle)]\right|\\
&\equiv Q_{2,1,k}(x, j_2) + Q_{2,2,k}(x, j_2).
\end{align*}

Therefore:
\[
|q_2(t, x, j_2, c_2)| \le \textcolor{red}{\sum_{k=1}^r |L_{C_2(j_2),k}|}(Q_{2,1,k}(x, j_2) + Q_{2,2,k}(x, j_2)) \le \textcolor{red}{rK}\max_{1 \le k \le r}(Q_{2,1,k}(x, j_2) + Q_{2,2,k}(x, j_2)),
\]
where we used $\|L\|_{\infty,1} \le \textcolor{red}{rK}$.

\paragraph*{Step 3: Decomposition of $q_{B,k}$ with low-rank structure.}

For the backpropagated signal difference:
\begin{align*}
|q_{B,k}(t, x)| &= \left|\frac{1}{n_2}\sum_{j_2=1}^{n_2} \textcolor{red}{L_{C_2(j_2),k}}\,\varphi_2'(\tilde{H}_2(t, j_2; x, \tilde{W}))\,\tilde{w}_2(t, j_2)\right.\\
&\quad\left. - \mathbb{E}_{C_2}[\textcolor{red}{L_{C_2,k}}\,\varphi_2'(H_2(t, C_2; x, W))\,w_2(t, C_2)]\right|.
\end{align*}

We decompose:
\begin{align*}
|q_{B,k}(t, x)| &\le \max_{j_2 \le n_2}\left|\textcolor{red}{L_{C_2(j_2),k}}\,\varphi_2'(\tilde{H}_2(t, j_2; x, \tilde{W}))\,\tilde{w}_2(t, j_2)\right.\\
&\quad\left. - \textcolor{red}{L_{C_2(j_2),k}}\,\varphi_2'(H_2(t, C_2(j_2); x, W))\,w_2(t, C_2(j_2))\right|\\
&\quad + \left|\frac{1}{n_2}\sum_{j_2=1}^{n_2} \textcolor{red}{L_{C_2(j_2),k}}\,\varphi_2'(H_2(t, C_2(j_2); x, W))\,w_2(t, C_2(j_2))\right.\\
&\quad\left. - \mathbb{E}_{C_2}[\textcolor{red}{L_{C_2,k}}\,\varphi_2'(H_2(t, C_2; x, W))\,w_2(t, C_2)]\right|\\
&\equiv Q_{B,1,k}(x) + Q_{B,2,k}(x).
\end{align*}

For $Q_{B,1,k}$, using Assumption~\ref{assump:regularity} and the fact that $\textcolor{red}{|L_{C_2(j_2),k}| \le K}$:
\begin{align*}
\mathbb{E}_Z[Q_{B,1,k}(X)] &\le K\max_{j_2 \le n_2}\left(|\tilde{w}_2(t, j_2) - w_2(t, C_2(j_2))|\right.\\
&\quad\left. + |w_2(t, C_2(j_2))|\,\mathbb{E}_Z[|\tilde{H}_2(t, j_2; X, \tilde{W}) - H_2(t, C_2(j_2); X, W)|]\right)\\
&\le K_t\left(\mathscr{D}_t(W, \tilde{W}) + \max_{j_2 \le n_2}\mathbb{E}_Z[|q_2(t, X, j_2, C_2(j_2))|]\right).
\end{align*}

\paragraph*{Step 4: Concentration bounds adapted to low-rank.}

For $Q_{2,1,k}$, we have:
\begin{align*}
\max_{j_2 \le n_2}\mathbb{E}_Z[Q_{2,1,k}(X, j_2)] &\le K\max_{j_1 \le n_1, 1 \le k \le r}|\tilde{w}_1(t, j_1, k) - w_1(t, C_1(j_1), k)|\\
&\le K_t\mathscr{D}_t(W, \tilde{W}).
\end{align*}

For $Q_{2,2,k}$, we apply concentration. Let us write:
\[
Z_{1,k}(x, c_1) = w_1(t, c_1, k)\,\varphi_1(\langle f_1(c_1), x\rangle).
\]
Recall that $\{C_1(j_1)\}_{j_1=1}^{n_1}$ are i.i.d. We have:
\[
\mathbb{E}[Z_{1,k}(X, C_1(j_1))|X] = \mathbb{E}_{C_1}[Z_{1,k}(X, C_1)],
\]
and $\{Z_{1,k}(X, C_1(j_1))\}_{j_1=1}^{n_1}$ are independent conditional on $X$. By Assumption~\ref{assump:regularity}, $|Z_{1,k}(X, C_1(j_1))| \le K_t$ almost surely. Then by concentration inequalities:
\[
\mathbb{P}\left(\mathbb{E}_Z[Q_{2,2,k}(X, j_2)] \ge K_t\gamma_{2,k}\right) \le \frac{1}{\gamma_{2,k}}\exp\left(-\frac{n_1\gamma_{2,k}^2}{K_t}\right).
\]

For $Q_{B,2,k}$, we note that almost surely:
\[
\textcolor{red}{|L_{C_2(j_2),k}|}\,\varphi_2'(H_2(t, C_2(j_2); x, W))\,w_2(t, C_2(j_2))| \le K\textcolor{red}{|L_{C_2(j_2),k}|} \le K^2,
\]
by Assumption~\ref{assump:regularity}. Then by concentration inequalities:
\[
\mathbb{P}\left(\mathbb{E}_Z[Q_{B,2,k}(X)] \ge K_t\gamma_{1,k}\right) \le \frac{1}{\gamma_{1,k}}\exp\left(-\frac{n_2\gamma_{1,k}^2}{K_t}\right).
\]

\paragraph*{Step 5: Combining bounds with union over $k$.}

\textbf{Taking a union bound over $k = 1, \ldots, r$ channels} (which is specific to the low-rank structure), we obtain that for any $\gamma_1, \gamma_2 > 0$ and $t \ge 0$, the event:
\begin{align*}
&\max\Bigg\{\max_{j_2 \le n_2}\mathbb{E}_Z[|q_2(t, X, j_2, C_2(j_2))|],\\
&\qquad\max_{j_1 \le n_1, 1 \le k \le r}\mathbb{E}_Z[|q_{B,k}(t, X)|]\Bigg\}\\
&\quad \le K_t\textcolor{red}{(1+rK)}(\mathscr{D}_t(W, \tilde{W}) + \gamma_1 + \gamma_2),
\end{align*}
holds with probability at least:
\[
1 - \frac{\mathbf{n_1r}}{\gamma_1}\exp\left(-\frac{n_2\gamma_1^2}{\mathbf{K_t\textcolor{red}{(1+rK)^2}}}\right) - \frac{\mathbf{n_2r}}{\gamma_2}\exp\left(-\frac{n_1\gamma_2^2}{\mathbf{K_t\textcolor{red}{(1+rK)^2}}}\right).
\]

\paragraph*{Step 6: Gronwall argument.}

Combining the bounds and taking a union bound over a discrete time grid $t \in \{0, \xi, 2\xi, \ldots, \lfloor T/\xi\rfloor\xi\}$ for some $\xi \in (0,1)$, we obtain:
\[
\max\left\{\max_{j_2 \le n_2}\left|\frac{\partial}{\partial t}\tilde{w}_2(t, j_2) - \frac{\partial}{\partial t}w_2(t, C_2(j_2))\right|,\right.
\]
\[
\left.\max_{j_1 \le n_1, 1 \le k \le r}\left|\frac{\partial}{\partial t}\tilde{w}_1(t, j_1, k) - \frac{\partial}{\partial t}w_1(t, C_1(j_1), k)\right|\right\}
\]
\[
\le \mathbf{K_T\textcolor{red}{(1+rK)}}(\mathscr{D}_t(W, \tilde{W}) + \gamma_1 + \gamma_2 + \xi), \qquad \forall t \in [0, T],
\]
with probability at least:
\[
1 - \frac{T+1}{\xi}\left[\frac{\mathbf{n_1r}}{\gamma_1}\exp\left(-\frac{n_2\gamma_1^2}{\mathbf{K_T\textcolor{red}{(1+rK)^2}}}\right) + \frac{\mathbf{n_2r}}{\gamma_2}\exp\left(-\frac{n_1\gamma_2^2}{\mathbf{K_T\textcolor{red}{(1+rK)^2}}}\right)\right].
\]

The above event implies:
\[
\mathscr{D}_t(W, \tilde{W}) \le \mathbf{K_T\textcolor{red}{(1+rK)}}\int_0^t(\mathscr{D}_s(W, \tilde{W}) + \gamma_1 + \gamma_2 + \xi)ds,
\]
and hence by Gronwall's lemma and the fact $\mathscr{D}_0(W, \tilde{W}) = 0$:
\[
\mathscr{D}_T(W, \tilde{W}) \le (\gamma_1 + \gamma_2 + \xi)\mathbf{e^{K_T\textcolor{red}{(1+rK)}}}.
\]

The result follows from choosing:
\[
\xi = \frac{1}{\sqrt{n_{\max}}}, \quad \gamma_1 = \frac{\mathbf{K_T\textcolor{red}{(1+rK)}}}{\sqrt{n_2}}\log^{1/2}\left(\frac{3(T+1)n_{\max}^2\mathbf{r}}{\delta} + e\right),
\]
\[
\gamma_2 = \frac{\mathbf{K_T\textcolor{red}{(1+rK)}}}{\sqrt{n_1}}\log^{1/2}\left(\frac{3(T+1)n_{\max}^2\mathbf{r}}{\delta} + e\right).
\]
\end{proof}

\subsection{Proof of Lemma~\ref{lem:gradient-lowrank}}

\begin{proof}
The proof follows the same structure as the full-width case \cite{nguyen2020global}, with careful adaptations for the low-rank structure. \textbf{The key difference is that we need to account for the $r$ channels in $w_1$ and the mixing matrix $L$ in all bounds.}

We consider $t \le T$ for a given terminal time $T \in \epsilon\mathbb{N}_{\ge 0}$. We reuse the notation $K_t$ from the proof of Lemma~\ref{lem:particle-lowrank}. Note that $K_t \le K_T$ for all $t \le T$. We also note that at initialization, $\mathscr{D}_0(\mathbf{W}, \tilde{W}) = 0$.

For brevity, let us define quantities that relate to the difference in the gradient updates between $\mathbf{W}$ and $\tilde{W}$:
\begin{align*}
q_2(k, z, \tilde{z}, j_2) &= d_L(z; \mathbf{W}(k))\,\varphi_2(\mathbf{H}_2(k, j_2; x, \mathbf{W}(k)))\\
&\quad - d_L(\tilde{z}; \tilde{W}(k\epsilon))\,\varphi_2(\tilde{H}_2(k\epsilon, j_2; \tilde{x}, \tilde{W}(k\epsilon))),\\
\textcolor{red}{r_2(k, z, j_2)} &= \xi_2(k\epsilon)d_L(z; \tilde{W}(k\epsilon))\,\varphi_2(\tilde{H}_2(k\epsilon, j_2; x, \tilde{W}(k\epsilon)))\\
&\quad - \xi_2(k\epsilon)\mathbb{E}_Z[d_L(Z; \tilde{W}(k\epsilon))\,\varphi_2(\tilde{H}_2(k\epsilon, j_2; X, \tilde{W}(k\epsilon)))],\\
q_1(k, z, \tilde{z}, j_1, k') &= d_L(z; \mathbf{W}(k))\,\varphi_1(\langle f_1(C_1(j_1)), x\rangle)\,\mathbf{B}_{k'}(k; x, \mathbf{W}(k))\\
&\quad - d_L(\tilde{z}; \tilde{W}(k\epsilon))\,\varphi_1(\langle f_1(C_1(j_1)), \tilde{x}\rangle)\,\tilde{B}_{k'}(k\epsilon; \tilde{x}, \tilde{W}(k\epsilon)),\\
\textcolor{red}{r_1(k, z, j_1, k')} &= \xi_1(k\epsilon)d_L(z; \tilde{W}(k\epsilon))\,\varphi_1(\langle f_1(C_1(j_1)), x\rangle)\,\tilde{B}_{k'}(k\epsilon; x, \tilde{W}(k\epsilon))\\
&\quad - \xi_1(k\epsilon)\mathbb{E}_Z[d_L(Z; \tilde{W}(k\epsilon))\,\varphi_1(\langle f_1(C_1(j_1)), X\rangle)\,\tilde{B}_{k'}(k\epsilon; X, \tilde{W}(k\epsilon))],
\end{align*}
where $\mathbf{B}_{k'}(k; x, \mathbf{W}) = \frac{1}{n_2}\sum_{j_2=1}^{n_2} L_{C_2(j_2),k'}\,\varphi_2'(\mathbf{H}_2(k, j_2; x, \mathbf{W}))\,\mathbf{w}_2(k, j_2)$.

By time-interpolation estimates (similar to Claim 1 in the full-width case) and Assumption~\ref{assump:regularity}, we have:
\begin{align*}
|\mathbf{w}_2(\lfloor t/\epsilon\rfloor, j_2) - \tilde{w}_2(t, j_2)| &\le K\max_{j_2 \le n_2}\left[Q_{2,1}(\lfloor t/\epsilon\rfloor, j_2) + Q_{2,2}(\lfloor t/\epsilon\rfloor, j_2)\right] + tK_t\epsilon,\\
|\mathbf{w}_1(\lfloor t/\epsilon\rfloor, j_1, k) - \tilde{w}_1(t, j_1, k)| &\le K\max_{j_1 \le n_1, 1 \le k \le r}\left[Q_{1,1}(\lfloor t/\epsilon\rfloor, j_1, k) + Q_{1,2}(\lfloor t/\epsilon\rfloor, j_1, k)\right] + tK_t\epsilon,
\end{align*}
where:
\begin{align*}
Q_{2,1}(\lfloor t/\epsilon\rfloor, j_2) &= \epsilon\sum_{\ell=0}^{\lfloor t/\epsilon\rfloor - 1}|q_2(\ell, z(\ell), z(\ell), j_2)|,\\
Q_{2,2}(\lfloor t/\epsilon\rfloor, j_2) &= \left|\epsilon\sum_{\ell=0}^{\lfloor t/\epsilon\rfloor - 1}\textcolor{red}{r_2(\ell, z(\ell), j_2)}\right|,\\
Q_{1,1}(\lfloor t/\epsilon\rfloor, j_1, k) &= \epsilon\sum_{\ell=0}^{\lfloor t/\epsilon\rfloor - 1}|q_1(\ell, z(\ell), z(\ell), j_1, k)|,\\
Q_{1,2}(\lfloor t/\epsilon\rfloor, j_1, k) &= \left|\epsilon\sum_{\ell=0}^{\lfloor t/\epsilon\rfloor - 1}\textcolor{red}{r_1(\ell, z(\ell), j_1, k)}\right|.
\end{align*}

\paragraph*{Bounding the terms:}

For $Q_{2,1}$, using Assumption~\ref{assump:regularity} and the low-rank structure:
\begin{align*}
|q_2(\ell, z(\ell), z(\ell), j_2)| &\le K|d_L(z(\ell); \mathbf{W}(\ell)) - d_L(z(\ell); \tilde{W}(\ell\epsilon))|\\
&\quad + K|d_L(z(\ell); \tilde{W}(\ell\epsilon))|\,|\mathbf{H}_2(\ell, j_2; x(\ell), \mathbf{W}(\ell)) - \tilde{H}_2(\ell\epsilon, j_2; x(\ell), \tilde{W}(\ell\epsilon))|\\
&\le K_t\mathscr{D}_{\ell\epsilon}(\tilde{W}, \mathbf{W}) + K_t\textcolor{red}{\sum_{k=1}^r |L_{C_2(j_2),k}|}\max_{j_1 \le n_1}|\mathbf{w}_1(\ell, j_1, k) - \tilde{w}_1(\ell\epsilon, j_1, k)|\\
&\le \mathbf{K_t\textcolor{red}{(1+rK)}}\mathscr{D}_{\ell\epsilon}(\tilde{W}, \mathbf{W}),
\end{align*}
which yields:
\[
\max_{j_2 \le n_2}Q_{2,1}(\lfloor t/\epsilon\rfloor, j_2) \le K_t\textcolor{red}{(1+rK)}\epsilon\sum_{\ell=0}^{\lfloor t/\epsilon\rfloor - 1}\mathscr{D}_{\ell\epsilon}(\tilde{W}, \mathbf{W}).
\]

For $Q_{1,1}$, similarly:
\begin{align*}
|q_1(\ell, z(\ell), z(\ell), j_1, k)| &\le K|d_L(z(\ell); \mathbf{W}(\ell)) - d_L(z(\ell); \tilde{W}(\ell\epsilon))|\,|\mathbf{B}_k(\ell; x(\ell), \mathbf{W}(\ell))|\\
&\quad + K|d_L(z(\ell); \tilde{W}(\ell\epsilon))|\,|\mathbf{B}_k(\ell; x(\ell), \mathbf{W}(\ell)) - \tilde{B}_k(\ell\epsilon; x(\ell), \tilde{W}(\ell\epsilon))|\\
&\le K_t\textcolor{red}{(1+rK)}\mathscr{D}_{\ell\epsilon}(\tilde{W}, \mathbf{W}),
\end{align*}
which yields:
\[
\max_{j_1 \le n_1, 1 \le k \le r}Q_{1,1}(\lfloor t/\epsilon\rfloor, j_1, k) \le K_t\textcolor{red}{(1+rK)}\epsilon\sum_{\ell=0}^{\lfloor t/\epsilon\rfloor - 1}\mathscr{D}_{\ell\epsilon}(\tilde{W}, \mathbf{W}).
\]

For $Q_{2,2}$ and $Q_{1,2}$, we use martingale concentration. The martingale differences are bounded: $\textcolor{red}{|r_2(\ell, z(\ell), j_2)| \le K_t}$ and $\textcolor{red}{|r_1(\ell, z(\ell), j_1, k)| \le K_t\textcolor{red}{(1+rK)}}$ almost surely by Assumption~\ref{assump:regularity} and the low-rank structure. Then by martingale concentration inequalities:
\[
\textcolor{red}{\mathbb{P}\left(\max_{j_2 \le n_2}\max_{\ell \in \{0,1,\ldots,T/\epsilon\}}Q_{2,2}(\ell, j_2) \ge \xi\right) \le 2n_2\exp\left(-\frac{\xi^2}{K_T(T+1)\epsilon}\right),}
\]
\[
\textcolor{red}{\mathbb{P}\left(\max_{j_1 \le n_1, 1 \le k \le r}\max_{\ell \in \{0,1,\ldots,T/\epsilon\}}Q_{1,2}(\ell, j_1, k) \ge \xi\right) \le \mathbf{2n_1r}\exp\left(-\frac{\xi^2}{\mathbf{K_T\textcolor{red}{(1+rK)^2}}(T+1)\epsilon}\right).}
\]

\paragraph*{Putting everything together:}

All the above results give us:
\[
\mathscr{D}_{\lfloor t/\epsilon\rfloor\epsilon}(\tilde{W}, \mathbf{W}) \le \mathbf{K_T\textcolor{red}{(1+rK)}}\epsilon\sum_{\ell=0}^{\lfloor t/\epsilon\rfloor - 1}\mathscr{D}_{\ell\epsilon}(\tilde{W}, \mathbf{W}) + \xi + TK_T\epsilon, \qquad \forall t \le T,
\]
which holds with probability at least:
\[
1 - \mathbf{2n_1r}\exp\left(-\frac{\xi^2}{\mathbf{K_T\textcolor{red}{(1+rK)^2}}(T+1)\epsilon}\right) - 2n_2\exp\left(-\frac{\xi^2}{K_T(T+1)\epsilon}\right).
\]

The above event implies, by Gronwall's lemma:
\[
\mathscr{D}_T(\tilde{W}, \mathbf{W}) \le (\xi + \epsilon)e^{K_T\textcolor{red}{(1+rK)}}.
\]

Choosing $\xi = \mathbf{K_T\textcolor{red}{(1+rK)}}\sqrt{(T+1)\epsilon\log(\mathbf{2n_1n_2r}/\delta)}$ completes the proof.
\end{proof}

